\documentclass{beamer}


\usepackage{tikz, tikz-cd}


\newcommand{\ol}{\overline}



\usetheme{madrid}

\title[Lagrangian fibrations]{The algebraic monodromy of Lagrangian fibrations}
\author{Edward Varvak}
\institute[UIC]{
    Department of Mathematics, Science, and Computer Science \\
    University of Illinois Chicago
}
\date{May 2026}


\begin{document}

\frame{\titlepage}

\begin{frame}
\frametitle{Holomorphic symplectic manifolds}

\begin{block}{Definition}
A {\it symplectic manifold} \((M, \omega)\) is a smooth manifold \(M\) equipped with a 
closed non-degenerate global 2-form \(\omega \in H^0\left(M, \bigwedge^2 T^* M\right)\).
\end{block}

Contracting with \(\omega\) induces an isomorphism 
\(\iota_{\omega}: TM \xrightarrow{\sim} T^* M\). \pause

\begin{examples}
\begin{itemize}
    \item<2-> \(\mathbb{R}^{2n}\) admits a natural symplectic form 
    \(\omega = \begin{pmatrix}
    0 & I_n \\ -I_n & 0
    \end{pmatrix}\). 
    \item<3-> Any K\"ahler manifold \(X\) is symplectic with respect to the K\"ahler form.
    \item<4-> Any projective manifold is symplectic with respect to the Fubini--Study metric.
\end{itemize}
\end{examples}
\pause[5]
Note that in the latter two examples, despite \(X\) being a complex manifold, 
the K\"ahler form \(\omega \sim dz_1 \wedge d\ol{z}_1 + \cdots + dz_n \wedge d\ol{z}_n\)
is not holomorphic.
\end{frame}




\begin{frame}{Holomorphic symplectic manifolds II}

\begin{block}{Definition}
A {\it holomorphic symplectic manifold} is a pair \((X, \sigma)\), where \(X\) is a 
K\"ahler manifold equipped with a closed non-degenerate 2-form 
\(\sigma \in H^0(X, \Omega_X^2)\).
\end{block}

\pause
Again, this gives an isomorphism 
\(\iota_\sigma: \mathcal{T}_X \xrightarrow{\sim} \Omega_X^1\).
\pause

\begin{examples}
    \begin{itemize}
        \item<3-> \(\mathbb{C}^{2n}\) admits a symplectic form \(\begin{pmatrix}
        0 & I_n \\ -I_n & 0
        \end{pmatrix}\).
        \item<4-> Complex tori \(\mathbb{C}^{2n}/\Lambda\), 
        \(\Lambda \simeq \mathbb{Z}^{4n}\), are symplectic.
        \item<5-> K3 surfaces \(S\), \(\pi_1(S) = \{1\}\), \(\Omega_S^2 \simeq \mathcal{O}_S\), 
        are symplectic.
        \item<6-> Any product of the spaces above is symplectic.
    \end{itemize}
\end{examples}

\end{frame}



\begin{frame}{The decomposition theorem}

We can generalize the K3 surface as follows:

\begin{block}{Definition}
    We say a compact holomorphic symplectic manifold \((X, \sigma)\) is 
    {\it irreducible} if it is simply connected and 
    \(\bigoplus_{p=0}^{\dim X} H^0(X, \Omega_X^{2p}) = \mathbb{C}[\sigma]\).
\end{block}
\pause

\begin{block}{Beauville--Bogomolov Decomposition Theorem}
    Let \((X, \sigma)\) be a compact holomorphic symplectic manifold. Then there exists 
    a finite cover \[\pi: T \times \prod_i X_i \to X,\]
    where \(T\) is a complex torus and the \(X_i\) are irreducible holomorphic 
    symplectic manifolds. 
\end{block}
\pause

Since complex tori are somewhat well understood, we will focus on understanding 
the irreducible holomorphic symplectic (IHS) \(X_i\).

\end{frame}




\begin{frame}

\begin{definition}
    Let \((X, \sigma)\) be a holomorphic symplectic manifold of dimension \(2n\), and 
    let \(Y \subseteq X\) be a subvariety. 
    \begin{enumerate}
        \item \(Y\) is {\it isotropic} if \(\sigma|_Y = 0\).
        \item \(Y\) is {\it Lagrangian} if \(Y\) is maximally isotropic, i.e. 
        \(\sigma|_Y = 0\) and \(\dim Y = n\).
    \end{enumerate}
\end{definition}

\pause
In general, we study symplectic manifolds by finding Lagrangian subvarieties.
\pause

\begin{theorem}[Greb--Lehn `14]
    Let \(f: X \to B\) be a proper surjective morphism, \(B\) a normal analytic space of 
    dimension not equal to \(0\) or \(2n\). Then the fibers of \(f\) are Lagrangian
    subvarieties and \(B\) is a projective variety of dimension \(n\).
\end{theorem}

\pause
For this reason, we refer to any such morphism \(f: X \to B\) as a 
{\it Lagrangian fibration}. Understanding the geometry of Lagrangian fibrations 
will help us understand the geometry of IHS manifolds.

\end{frame}



\begin{frame}{The structure of the fibration}

Let \(f^\circ: X^\circ \to B^\circ\) be the smooth fibers of \(f: X \to B\). 
We need to understand 
\begin{itemize}
    \item the base \(B\),
    \item the general fiber \(X_b\),
    \item the action of \(\pi_1(B^\circ, b)\) on \(X_b\).
\end{itemize}
\pause

\begin{theorem}
    \begin{itemize}
        \item (Matsushita `99) The smooth fibers \(X_b\) are Abelian varieties. \pause
        \item (Hwang `08) If \(B\) is smooth, \(B \simeq \mathbb{P}^n\).
    \end{itemize}
\end{theorem}

\pause
By Hodge theory, an Abelian variety \(A\) can be recovered from its cohomology:
\begin{equation*}
    A \simeq H^{1,0}(A) / (H^1(A, \mathbb{Z}))^\vee.
\end{equation*}
We will study the monodromy representation 
\[\rho: \pi_1(B^\circ, b) \to \mathrm{Aut}(H^1(X_b, \mathbb{Z})).\]


\end{frame}


\begin{frame}{Algebraic monodromy groups}

\begin{Definition}
The {\it full algebraic monodromy group} \(M\) of \(\rho\) is the Zariski closure of 
\(\rho(\pi_1(B^\circ, b)) \subseteq \mathrm{Aut}(H^1(X_b, \mathbb{Q})) = 
\mathrm{Sp}_{2n, \mathbb{Q}}\). Its identity component 
\(M^\circ \subseteq M\) is the {\it (connected) algebraic monodromy group}.
\end{Definition}
\pause

When \(X\) is projective, we can classify the possible representations of 
the algebraic monodromy group.

\begin{Theorem}[V `26+]
Let \(X\) be a \(2n\)-dimensional projective IHS manifold admitting a Lagrangian fibration 
\(f: X \to B\). The algebraic monodromy group \(M^\circ\) associated to \(f\) is 
one of the following groups: \pause
\begin{enumerate}
    \item \(M^\circ = \mathrm{Sp}_{2n, \mathbb Q}\), acting on \(H^1(X_b, \mathbb Q)\) 
    via the standard representation; \pause
    \item \(M^\circ = (\mathrm{SL}_{2, \mathbb Q})^n\), acting on \(H^1(X_b, \mathbb Q)\)
    diagonally; \pause
    \item \(M^\circ = \{1\}\), acting trivially on \(H^1(X_b, \mathbb{Q})\).
\end{enumerate}
\end{Theorem}

\end{frame}



\begin{frame}{Geometric consequences}

In each of the cases above, we extract information regarding the geometry 
of \(X\) and \(f: X \to B\). \pause

\begin{Theorem}[V `26+]
\begin{enumerate}
    \item When \(M^\circ = \mathrm{Sp}_{2n, \mathbb Q}\), the very general fiber 
    of \(f\) is a simple Abelian variety with no extra endomorphisms. \pause
    \item When \(M^\circ = (\mathrm{SL_{2, \mathbb Q}})^n\), there exists an 
    elliptically fibered K3 surface \(g: S \to \mathbb{P}^1\) such that 
    \(X \to B\) is birational to \(S^n/\mathfrak{S}_n \to \mathbb{P}^n\).
    In particular, the fibers \(X_b\) are isogenous to products of elliptic curves.
\end{enumerate}
\end{Theorem}
\pause

The case (3) is also somewhat well understood based on previous work:

\begin{Theorem}[Kim--Laza--Martin `25, V `26]
    When \(M^\circ = \{1\}\), there exists an elliptic curve \(E\) 
    for which the smooth fibers \(X_b\) are isogenous to \(E^n\).
\end{Theorem}

\end{frame}



\begin{frame}[fragile]{Proof of the main theorem.}

{\it Proof idea.} 
We have the following isomorphism of bundles:
\begin{equation*}
\begin{tikzcd}
0 \arrow{r} & \mathcal{T}_{X^\circ/B^\circ} \arrow{r} \arrow[equals]{d} & \mathcal{T}_{X^\circ} \arrow{r} \arrow[equals]{d} & f^*\mathcal{T}_{B^\circ} \arrow{r} \arrow[equals]{d} & 0 \\
0 \arrow{r} & f^*\Omega_{B^\circ}^1 \arrow{r} & \Omega_{X^\circ}^1 \arrow{r} & \Omega_{X^\circ/B^\circ}^1 \arrow{r} & 0
\end{tikzcd}
\end{equation*}
The fibers of \(f\) are given by quotienting \(f_*\Omega^1_{X^\circ/B^\circ, b}\) 
by a full rank sublattice. In particular, the structure of the fibers of 
\(f\) comes from \(\mathcal{T}_B\). \pause \newline

There exists a cover \(\pi: B' \to B\) for which 
\(\mathcal{T}_{B'} = \mathcal{F}_1 \oplus \cdots \oplus \mathcal{F}_k\) splits 
into transverse foliations \(\mathcal{F}_i\), which are permuted by the 
covering group \(\mathfrak{S}_k\). \pause
The general leaves of \(\mathcal{F}_i\) 
are algebraic, yielding a generically finite rational map 
\begin{equation*}
    p: B' \dasharrow Z^k.
\end{equation*}

\end{frame}



\begin{frame}[fragile]{Proof of main theorem.}

{\it Proof idea (continued).} Now quotient \(p: B' \dasharrow Z^k\) by \(\mathfrak{S}_k\).
We get a rational map \(X \xrightarrow{f} B \dasharrow Z^k/\mathfrak{S}_k\).
By a criterion of Sacc\`a, we can construct another IHS manifold \(\widetilde{X}\) 
birational to \(X\) with a morphism \(\varphi: \widetilde{X} \to Z^k/\mathfrak{S}_k\)
extending the composition. \pause \newline

By a theorem of Matsushita, \(Z^k/\mathfrak{S}_k\) has the same cohomology as 
\(\mathbb{P}^n\). This forces 
\begin{equation*}
    \mathrm{IH}^*(Z, \mathbb{Q}) = H^*(\mathbb{P}^1, \mathbb{Q}),
\end{equation*}
so that \(Z = \mathbb{P}^1\), and \(k = n\). \pause
Now we may construct a K3 surface \(S\) and a family of elliptic curves 
\(g: S \to \mathbb{P}^1\) so that 
\begin{equation*}
\begin{tikzcd}
S^n/\mathfrak{S}_n \arrow[dashed]{rr}{\sim} \arrow{dr} && \widetilde{X} \arrow{dl}{\varphi} \\
& \mathbb{P}^n &
\end{tikzcd}
\end{equation*}
\qed

\end{frame}


\begin{frame}{Questions}
    
\end{frame}


\end{document}